\bigskip

\begin{theorem} \label{an:thm:isolated-points}
  Suppose $(X, d)$ is a metric space and $E \subseteq X$. Then
  \[
    \forall p \in X, p \in E \setminus E' \iff p \text{ is an isolated point of } E
  \]
\end{theorem}

\begin{proof}
  \textbf{($\implies$)} Suppose $p \in E \setminus E'$. Then
  \[
    \begin{array}{c}
      p \notin E'                                               \\
      \Downarrow                                                \\
      \exists r \in \R^{+}, E \cap \widetilde{N}_{r}(p) = \vnot \\
      \Downarrow                                                \\
      \paren*{E \cap N_{r}(p)} \setminus \set*{p} = \vnot
    \end{array}
  \]
  Therefore,
  \[
    \begin{array}{c}
      p \in E \land p \in N_{r}(p) \\
      \Downarrow                   \\
      E \cap N_{r}(p) = \set*{p}   \\
      \Downarrow                   \\
      p \text{ is an isolated point of } E
    \end{array}
  \]
  \textbf{($\impliedby$)} Suppose $p$ is an isolated point of $E$. Then
  \[
    \exists r \in \R^{+}, E \cap N_{r}(p) = \set*{p}
  \]
  Therefore,
  \[
    \begin{array}{c}
      p \in E                             \\
      \Downarrow                          \\
      E \cap \widetilde{N}_{r}(p) = \vnot \\
      \Downarrow                          \\
      p \notin E'                         \\
      \Downarrow                          \\
      p \in E \setminus E'
    \end{array}
  \]
\end{proof}

\bigskip

\begin{definition} \label{an:def:boundary}
  Suppose $(X, d)$ is a metric space and $E \subseteq X$. The \textbf{boundary} $\partial E$ of $E$ is defined as
  \[
    \partial E \coloneqq \overline{E} \cap \overline{E^{c}}
  \]
\end{definition}

\bigskip

\begin{theorem}
  Suppose $(X, d)$ is a metric space and $E \subseteq X$. Then
  \[
    \partial E = \overline{E} \setminus E^{\circ}
  \]
\end{theorem}

\begin{proof}
  We will show
  \[
    \begin{array}{c}
      \paren*{E^{\circ}}^{c} = \overline{E^{c}} \\
      \Downarrow                                \\
      \overline{E} \cap \overline{E^{c}} = \overline{E} \setminus E^{\circ}
    \end{array}
  \]
  Suppose $p \in X$. Then
  \begin{align*}
    p \in \paren*{E^{\circ}}^{c} & \iff p \notin E^{\circ}                                   \\
                                 & \iff \forall r \in \R^{+}, N_{r}(p) \not\subseteq E       \\
                                 & \iff \forall r \in \R^{+}, N_{r}(p) \cap E^{c} \neq \vnot \\
                                 & \iff p \in \overline{E^{c}}
  \end{align*}
\end{proof}

\bigskip

\begin{lemma} \label{an:lem:deleted-neighborhood-in-rn-is-non-empty}
  Suppose $n \in \N$. Then
  \[
    \forall \bm{p} \in \R^{n}, \forall r \in \R^{+}, \widetilde{N}_r(\bm{p}) \neq \vnot
  \]
\end{lemma}

\begin{proof}
  Suppose $\bm{p} \in \R^{n}$ and $r \in \R^{+}$. Then
  \[
    d\paren*{\bm{p} + \frac{r}{2} \bm{e_{1}}, \bm{p}} = \norm*{\paren*{\bm{p} + \frac{r}{2} \bm{e_{1}}}
      - \bm{p}} = \norm*{\frac{r}{2} \bm{e_{1}}} = \frac{r}{2} \norm*{\bm{e_{1}}} = \frac{r}{2} < r
  \]
  where $\bm{e_{1}} = (1, 0, 0, \dots, 0) \in \R^{n}$. Therefore,
  \[
    \begin{array}{c}
      \bm{p} \neq \paren*{\bm{p} + \frac{r}{2} \bm{e_{1}}} \in \widetilde{N}_r(\bm{p}) \\
      \Downarrow                                                                       \\
      \widetilde{N}_r(\bm{p}) \neq \vnot
    \end{array}
  \]
\end{proof}

\bigskip

\begin{theorem} \label{an:thm:properties-of-boundary-in-rn}
  Suppose $n \in \N$ and $E \subseteq \R^{n}$. Then
  \[
    \partial E = \paren*{E \setminus E'} \cup \paren*{E^{c} \setminus (E^{c})'} \cup \paren*{E' \cap (E^{c})'}
  \]
  This states that $\partial E$ consists of
  \begin{itemize}
    \item The set of all isolated points of $E$
    \item The set of all isolated points of $E^{c}$
    \item The set of all limit points of both $E$ and $E^{c}$
  \end{itemize}
  where the three sets are disjoint.
\end{theorem}

\begin{proof}
  By definition,
  \begin{align*}
    \partial E & = \overline{E} \cap \overline{E^{c}}                   \\
               & = \paren*{E \cup E'} \cap \paren*{E^{c} \cup (E^{c})'} \\
               & = \paren*{E \cap E^{c}} \cup \paren*{E \cap (E^{c})'}
    \cup \paren*{E' \cap E^{c}} \cup \paren*{E' \cap (E^{c})'}
  \end{align*}
  Consider
  \begin{itemize}
    \item $E = \paren*{E \setminus E'} \cup \paren*{E \cap E'}$
    \item $E^{c} = \paren*{E^{c} \setminus (E^{c})'} \cup \paren*{E^{c} \cap (E^{c})'}$
  \end{itemize}
  Then
  \begin{itemize}
    \item $E \cap (E^{c})' = \paren*{\paren*{E \setminus E'} \cap (E^{c})'}
            \cup \paren*{\paren*{E \cap E'} \cap (E^{c})'}$
    \item $E' \cap E^{c} = \paren*{E' \cap \paren*{E^{c} \setminus (E^{c})'}}
            \cup \paren*{E' \cap \paren*{E^{c} \cap (E^{c})'}}$
  \end{itemize}
  Consider
  \begin{itemize}
    \item $\paren*{\paren*{E \cap E'} \cap (E^{c})'} \subseteq E' \cap (E^{c})'$
    \item $\paren*{E' \cap \paren*{E^{c} \cap (E^{c})'}} \subseteq E' \cap (E^{c})'$
  \end{itemize}
  Then
  \begin{align*}
    \partial E & = \vnot \cup \paren*{\paren*{E \setminus E'} \cap (E^{c})'}
    \cup \paren*{E' \cap \paren*{E^{c} \setminus (E^{c})'}} \cup \paren*{E' \cap (E^{c})'} \\
               & = \paren*{\paren*{E \setminus E'} \cap (E^{c})'}
    \cup \paren*{E' \cap \paren*{E^{c} \setminus (E^{c})'}} \cup \paren*{E' \cap (E^{c})'}
  \end{align*}
  Also,
  \[
    \begin{array}{c}
      \forall \bm{p} \in E \setminus E', \exists r \in \R^{+}, E \cap \widetilde{N}_{r}(\bm{p}) = \vnot \\
      \Downarrow                                                                                        \\
      \widetilde{N}_{r}(\bm{p}) \subseteq E^{c}
    \end{array}
  \]
  By \cref{an:lem:deleted-neighborhood-in-rn-is-non-empty},
  \[
    \forall s \in \R^{+}, \widetilde{N}_{s}(\bm{p}) \cap \widetilde{N}_{r}(\bm{p}) =
    \widetilde{N}_{\min(s,r)}(\bm{p}) \neq \vnot
  \]
  Then
  \[
    \begin{array}{c}
      \forall s \in \R^{+}, \widetilde{N}_{s}(\bm{p}) \cap \widetilde{N}_{r}(\bm{p}) \neq \vnot \\
      \Downarrow                                                                                \\
      \widetilde{N}_{s}(\bm{p}) \cap \paren*{\widetilde{N}_{r}(\bm{p}) \cap E^{c}} \neq \vnot   \\
      \Downarrow                                                                                \\
      \vnot \neq \paren*{\widetilde{N}_{s}(\bm{p}) \cap E^{c}} \cap \widetilde{N}_{r}(\bm{p})
      \subseteq \widetilde{N}_{s}(\bm{p}) \cap E^{c}                                            \\
      \Downarrow                                                                                \\
      E^{c} \cap \widetilde{N}_{s}(\bm{p}) \neq \vnot                                           \\
      \Downarrow                                                                                \\
      \bm{p} \in (E^{c})'
    \end{array}
  \]
  Then
  \[
    \begin{array}{c}
      E \setminus E' \subseteq (E^{c})' \\
      \Downarrow                        \\
      \paren*{E \setminus E'} \cap (E^{c})' = E \setminus E'
    \end{array}
  \]
  Similarly,
  \[
    \begin{array}{c}
      E^{c} \setminus (E^{c})' \subseteq E' \\
      \Downarrow                            \\
      E' \cap \paren*{E^{c} \setminus (E^{c})'} = E^{c} \setminus (E^{c})'
    \end{array}
  \]
  Therefore,
  \[
    \partial E = \paren*{E \setminus E'} \cup \paren*{E^{c} \setminus (E^{c})'} \cup \paren*{E' \cap (E^{c})'}
  \]
\end{proof}
